\begin{lemma}[Stable terminal colorings are mod colorings]
Let (T) be a positive-rung extended Tashkinov tree generated by
\[
 \mathcal T=\{(T_i,\varphi_{i-1},S_{i-1},F_{i-1},\Theta_{i-1}):
                   1\leq i\leq n+1\},
\]
and suppose that (T) has the maximum property.  Keep the graph (G),
palette size (k), and distinguished uncolored edge (e_0) fixed, but
allow the certified input coloring to vary arbitrarily.  Assume explicitly
the following six assertions for every positive-rung ETT (T') on these
fixed data, with arbitrary certified input coloring \(\varphi'_0\), maximum
property, and rung number (s<n), writing \(\varphi'_s\) for its generating
coloring and \(\mathcal T'\) for its certified series:
\begin{enumerate}
\item (V(T')) is elementary with respect to \(\varphi'_s\).
\item Under every proper \((T'_{s+1},D'_s,\varphi'_s)\)-stable coloring
leaving precisely (e_0) uncolored, (T'_{s+1}) is interchangeable for
any pair of colors of which at least one is missing on (T'_{s+1}).
\item For each PE slot (a\leq s) with supporting vertex (u_a) and
(m(u_a)=b), every proper \((T'_a,D'_a,\varphi'_a)\)-stable coloring
leaving precisely (e_0) uncolored is both
\((T'(u_a)-u_a,D'_{b-1},\varphi'_{b-1})\)-stable and
\((T'_{b-1},D'_{b-1},\varphi'_{b-1})\)-stable.  Moreover, the stored
color pairs at two distinct PE slots with distinct supporting vertices
and the same first-appearance index are disjoint.
\item If the terminal slot is PE with supporting vertex (u_s) and
stored colors \(\gamma'_s,\delta'_s\), then under every proper
\((T'_s,D'_s,\varphi'_s)\)-stable coloring leaving precisely (e_0)
uncolored, the \((\gamma'_s,\delta'_s)\)-component through (u_s)
contains exactly the one vertex (u_s) of (T'_s).
\item If \(\tau\) is any proper \((T'_s,D'_s,\varphi'_s)\)-stable
coloring leaving precisely (e_0) uncolored, (d) is defective on
(T'_s) under \(\tau\), and (x) is a nonminimum vertex of
\(I[\partial_{\tau,d}(T'_s)]\), then
(x\preceq_{T'}u_a) for every PE supporting vertex or SE extension
vertex (u_a) with (m(x)\leq a\leq s).
\item Every proper \((T'_s,D'_s,\varphi'_s)\)-stable coloring leaving
precisely (e_0) uncolored is a \(\varphi'_s\bmod T'_s\) coloring.
\end{enumerate}

Then every proper \((T_n,D_n,\varphi_n)\)-stable coloring \(\sigma_n\)
leaving precisely (e_0) uncolored is a
\(\varphi_n\bmod T_n\) coloring in the sense of
\(\noderef{ModColoring}\).  In particular, the conclusion is universal
over every eligible ordered tree-sequence strictly extending (T_n),
and not merely over one selected closure.
\end{lemma}
\begin{proof}
We first obtain the current-rung assertions used below, rather than adding
them as hypotheses.  Specialize the displayed lower-rung package to the
original input coloring of (T).  This is exactly the lower premise of
\noderef{ExtensionRules}, so current assertion (iii) holds.  Keep the
arbitrary-input quantifier in the package and apply
\noderef{ExtensionBase}; this gives current assertion (iv), the exact
one-vertex intersection for the terminal PE component.  Applying
\noderef{Basic2} to the same generalized package gives current assertion
(v), including the order conclusion for the chosen vertex at every PE or
SE slot.  Thus all three current facts are consequences of the stated
lower package and the given \noderef{MaximumProperty}.

Fix a proper \((T_n,D_n,\varphi_n)\)-stable coloring \(\sigma_n\) leaving
exactly (e_0) uncolored.  The first two clauses of
\noderef{ModColoring} are already hypotheses.  For its universal clause,
fix an arbitrary tree-sequence (T^*) such that (T_n) is an ordered
initial segment of (T^*), the length inequality is strict, and the
appropriate eligibility alternative holds.  Thus, in the PE case, an
exact \(\sigma_n\)-\noderef{TashkinovAugmentationSequence} runs from
(T_n) to (T^*).  In the SE or RE case there are an individual edge
identity (f_n), the exact equality (F_n=\{f_n\}), and an exact seed
(A=T_n+f_n) in the sense of \noderef{TreeSequenceAdjoinEdge}, followed
by such an augmentation sequence from (A) to (T^*).

Apply \noderef{TashkinovClosureExists} to (T^*) under \(\sigma_n\), and
call the resulting closure (C).  Its
\noderef{TashkinovClosure} certificate gives an augmentation sequence
from (T^*) to (C) and closedness of (V(C)).  Each augmentation
retains the old display as the exact ordered initial segment by
\noderef{TashkinovAugmentation}; hence composition gives
\[
                         T_n\prec T^*\preccurlyeq C.          \tag{1}
\]
In the PE alternative, transitivity of
\noderef{TashkinovAugmentationSequence} composes the route from (T_n)
through (T^*) with the closure route.  In the SE/RE alternative it
composes the route from the same exact one-edge seed (A), retaining the
identity (f_n).  We shall make (C), not (T^*), the last closure in a
replacement series, while making (T^*) the ETT certified between its
last two ladder trees.  This is what proves the universal statement for
the originally arbitrary (T^*).

We next reconstruct the replacement series through the preceding rung.
Inspect the terminal positive certificate in
\noderef{TashkinovSeries} and \noderef{TashkinovSeriesStep}.  Define
\(\sigma_{n-1}\) as follows:
\[
 \sigma_{n-1}=\sigma_n\quad(\Theta_n=\mathrm{RE}\text{ or }\mathrm{SE}),
 \qquad
 \sigma_{n-1}=\sigma_n/P_{v_n}(\gamma_n,\delta_n,\sigma_n)
       \quad(\Theta_n=\mathrm{PE}).                           \tag{2}
\]
In an RE, the constructor says \(\varphi_n=\varphi_{n-1}\).  The first
inclusion of \noderef{HkuMissingStoredMonotone}, together with the two
clauses of the assumed \noderef{StableColoring}, therefore restricts the
control set from (D_n) to (D_{n-1}) and proves that
\(\sigma_{n-1}\) is \((T_n,D_{n-1},\varphi_{n-1})\)-stable.  In an SE,
the constructor identifies \(\varphi_n\) with the selected coloring
\(\pi_{n-1}\), and its \noderef{MaximumDefectiveChoice} says that
\(\pi_{n-1}\) is \((T_n,D_{n-1},\varphi_{n-1})\)-stable.  Since
\(\delta_n\in S_n\), the definition \noderef{TashkinovStoredColors}
puts \(\delta_n\) in
\(\overline\varphi_n(T_n)\cup D_n\); terminal stability therefore also
restricts to control (D_{n-1}\cup\{\delta_n\}) relative to
\(\pi_{n-1}\).  Transitivity from \noderef{sc2} gives the same preceding
stability conclusion.

In a PE, current assertion (iv), obtained from
\noderef{ExtensionBase}, says exactly that the switched
\((\gamma_n,\delta_n)\)-component meets (V(T_n)) only in (v_n).
All hypotheses of \noderef{ParallelExtensionUndoStable} are therefore
satisfied.  That lemma proves directly that the second coloring in (2)
is proper, leaves exactly (e_0) uncolored, and is
\((T_n,D_{n-1},\varphi_{n-1})\)-stable.  It also supplies the selected
coloring, the second stable coloring, and all enlarged-control stability
relations that will be used to rebuild the terminal PE constructor.
Consequently, in all three cases,
\[
 \sigma_{n-1}\text{ is proper, leaves exactly }e_0\text{ uncolored, and is }
 (T_n,D_{n-1},\varphi_{n-1})\text{-stable}.                    \tag{3}
\]

We require the preceding series to end at the exact ordered tree (T_n).
The previous constructor of \noderef{TashkinovSeriesStep} presents (T_n)
as a closure of (T_{n-1}) in the PE alternative, or as a closure of the
exact seed (T_{n-1}+f_{n-1}) in the SE/RE alternative, under
\(\varphi_{n-1}\).  We copy this closure edge by edge under
\(\sigma_{n-1}\).  At a copied augmentation, its new edge identity is a
displayed edge of (T_n), and its old missing-color witness is a vertex
of (T_n).  The vertex clause of (3) preserves that missing color, while
\noderef{HkuTreeColorsControlled} and the edge clause of (3) preserve the
color of the same exact displayed edge identity.  Hence the defining
clauses of \noderef{TashkinovAugmentation} hold for every copied step.
For closedness, a \(\sigma_{n-1}\)-colored boundary identity whose color
were missing on (T_n) would have the same missing witness under
\(\varphi_{n-1}\); the reverse edge criterion \noderef{sc1} would give
the same color to that exact identity under \(\varphi_{n-1}\), contrary
to the old closure.  Thus (T_n) is a closure of the same exact seed
under \(\sigma_{n-1}\).                                      \tag{4}

Suppose first that (n>1).  Restrict the original certified series
through rung (n-1), and use (T_n) as the ETT displayed between its
last two stages.  The strict ordered-extension clauses of
\noderef{TashkinovSeriesStep} give an
\noderef{ExtendedTashkinovTree} (B) with (n-1) rungs and generating
coloring \(\varphi_{n-1}\).  The initial maximum clause and every
stage-bound below (n-1) are literally restrictions of the corresponding
clauses of the given \noderef{MaximumProperty}, so (B) has the maximum
property.  Apply lower assertion (vi) to (B) and (3).  It gives
\[
                         \noderef{ModColoring}(B,\sigma_{n-1}). \tag{5}
\]
The copied closure (4) makes the exact tree (T_n) an eligible strict
extension for the universal clause of (5), with the same PE or exact
SE/RE seed used by the preceding original constructor.  That clause
supplies an actual \noderef{CorrespondingExtendedTashkinovTree}
certificate on (T_n), not merely an existential closure.  Retain its
replacement input, all replacement colorings, and all constructors
through rung (n-1).  If its last closure tree is larger than (T_n),
replace that last closure witness by the exact closure (4), and rebuild
the terminal ordered-extension and strictness fields from its augmentation
sequence and the old strict seed.  These are the only other terminal
fields depending on the chosen closure.  This produces a certified
replacement prefix whose last tree is
exactly (T_n), whose last coloring is \(\sigma_{n-1}\), and whose
ordered ladder trees, raw stored sets, exact connecting-edge sets, and
kinds agree with the original through slot (n-1).  Applying
\noderef{MaximumPropertyTransport} before this harmless last-closure
normalization gives maximum property for the replacement ETT; that
property only uses ladder stages through rung (n-1), so the normalized
prefix retains it.

When (n=1), there is no positive lower ETT to which (vi) could apply.
Instead retain the original one-edge seed (T_0), take
\(\sigma_0=\sigma_{n-1}\) as replacement input, and use (4) for the
initial closure (T_0\to T_1).  The graph, chromatic-index, and critical
clauses of the original \(\noderef{KTriple}\) do not depend on the
coloring; properness and exact unique-uncoloredness from (3) provide its
coloring clauses for \(\sigma_0\).  Thus the initial constructor of
\noderef{TashkinovSeriesStep} gives the required certified replacement
initial tuple ending at the exact (T_1=T_n).  Notice that
\noderef{TechnicalTheoremBase} proves only rung-zero elementarity and
interchangeability; no mod-coloring conclusion has been inferred from it.
This is the required direct reconstruction at the one-rung boundary.

It remains to append the terminal positive constructor.  We do so in the
three possible cases, always using the closure (C) chosen after the
arbitrary (T^*).

\textbf{Terminal RE.}
The old RE constructor contains an actual
\noderef{RevisitingExtensionCertificate}: an earlier positive PE slot
(h\leq n-1), the exact colors \(\gamma_h,\delta_h\), a full incidence
cycle, the exact connecting edge (f_n) and its two endpoints, and a
return path from the edge-node of (f_n) to (T_h), every graph vertex
of that path lying in (T_n).  The replacement prefix agrees on the
earlier kind, gamma, delta, raw stored set, and exact connecting edge.
Every edge identity on the return path is incident with (T_n), and its
color is one of \(\gamma_h,\delta_h\).  Membership of those colors in
(S_h), unfolded through \noderef{TashkinovStoredColors}, puts each in
\(\overline\varphi_{n-1}(T_n)\cup D_{n-1}\).  Equation (3) therefore
fixes every one of those identities under \(\sigma_{n-1}\), while the
endpoint and path-incidence equations are purely structural.  The same
exact return path hence lies in one
\((\gamma_h,\delta_h)\)-component under \(\sigma_{n-1}\), meets (T_h),
contains (f_n), and reaches its recorded endpoint outside (T_n).

We verify that this component is a cycle, which is the nonformal part of
the RE reconstruction.  In the replacement prefix, the earlier PE
constructor switches its exit component and makes \(\delta_h\) missing
at its supporting vertex.  The intervening RE constructors keep the
coloring unchanged; equivalently, choosing the last earlier PE belonging
to the retained RE chain gives a nontrivial
\((\gamma_h,\delta_h)\)-path component meeting (T_n).  Its intersection
with (T_h) is only the supporting vertex by the lower instance of
assertion (iv), whereas the retained return path meets (T_h) at a
vertex of the old cycle, distinct from that supporting vertex.  Thus the
two components are distinct.  The lower instance of assertion (ii),
applied to the replacement prefix (whose maximum property was just
established), and the definition \noderef{TreeInterchangeable} allow at
most one nontrivial path component for these colors to meet (T_n).
Therefore the component containing the return path is not a path.  Every
node has degree at most two by
\noderef{TwoColorIncidenceDegreeAtMostTwo}; it is finite, connected, and
contains an edge-node, so it has a cyclic enumeration and supplies a full
\noderef{TwoColorCycle}.  At (n=1) this case is impossible, because the
positive earlier index required by
\noderef{RevisitingExtensionCertificate} would have to be at most zero.

Use this cycle, the unchanged edge identity (f_n), its exact endpoint
equation, and the retained return path to construct an actual
\noderef{RevisitingExtensionCertificate} for the replacement prefix.
Every field has now been checked: earlier PE data and colors agree,
the cycle meets (T_h) and leaves (T_n), (f_n) lies both on the cycle
and in the exact boundary, and all internal return-path vertices lie in
(T_n).  The RE constructor of \noderef{TashkinovSeriesStep} therefore
applies.  It keeps \(\sigma_n=\sigma_{n-1}\), stores the same exact
sets and kind, first adjoins the same edge (f_n), and uses the composed
closure route from that seed through (T^*) to (C).

\textbf{Terminal SE or PE.}
Let the old maximum-defective data be
\((\pi_{n-1},v_n,f_n,u_n,\delta_n)\).  In the SE case (2) gives
\(\sigma_{n-1}=\sigma_n\); in the PE case the complete stability
relations between \(\sigma_{n-1}\), the second stable coloring, and
\(\pi_{n-1}\) are supplied by
\noderef{ParallelExtensionUndoStable}.  In both cases
\(\sigma_{n-1}\) is stable relative to \(\pi_{n-1}\) for
(D_{n-1}\cup\{\delta_n\}).  The vertex clause preserves all missing
sets on (T_n).  The edge clause and its symmetric form
\noderef{sc1} preserve, by identity, the complete
\(\delta_n\)-colored boundary set; in particular the same individual
edge (f_n) remains defective, incident with (v_n), has outside
endpoint (u_n), and has color \(\delta_n\).

We now rebuild \(\noderef{MaximumDefectiveChoice}\), including its
quantified fields.  Choose \(\sigma_{n-1}\), the same vertex (v_n),
the same edge (f_n), and the same other endpoint and color.  Properness,
unique-uncoloredness, defective incidence, and all exact equations were
just established.  If \(\rho\) is any proper coloring leaving (e_0)
uncolored and stable relative to \(\sigma_{n-1}\) for (D_{n-1}), then
\noderef{sc2} composes that stability with (3), so \(\rho\) is admissible
for the old maximum-defective choice.  Its universal maximality gives
(w=v_n) or (w\prec_{T_n}v_n) for every defective vertex (w).
This proves the new universal maximality field.  It is also consistent
with current assertion (v) from \noderef{Basic2}: after the same exact
boundary transport, every nonminimum defective vertex lies no later than
the terminal PE support or SE extension vertex, while the minimum case is
immediate from the tree order.

For SE, take an arbitrary coloring stable relative to
\(\sigma_{n-1}\) with enlarged control
(D_{n-1}\cup\{\delta_n\}).  Stability equivalence from \noderef{sc2}
transports it to the old selected coloring \(\pi_{n-1}\), so the old
universal series-alternative field gives the same disjointness of the
missing colors at (u_n) and on (T_n).  For PE, transport the old
existential counterexample in the reverse direction; missing sets are
unchanged vertexwise, so its failure of disjointness is unchanged.  Thus
the exact SE universal field or PE failure field is rebuilt, not merely
asserted.

The history part of \noderef{ParallelExtensionColor} depends only on the
earlier kinds, chosen vertices, and delta fields, all of which agree in
the replacement prefix.  Its remaining clause is that \(\gamma_n\) is
missing at (v_n); this is transported by the vertex clause of the
enlarged-control stability.  Hence the same \(\gamma_n\) obeys the PE
rule.  In the PE case choose \(\sigma_{n-1}\) itself as the second stable
coloring.  Reflexivity in \noderef{sc2} gives the required enlarged-control
stability.  The component used in (2) has the same support before and
after its color transposition.  Current assertion (iv) says that it meets
(T_n) only at (v_n); properness and the exact (f_n)-incidence make
(v_n) an endpoint, and \noderef{TwoColorIncidenceDegreeAtMostTwo}
then makes the finite component an exit path with its other endpoint
outside (T_n).  This is the required PE exit witness.  The involutive
definition \noderef{TwoColorComponentSwap} gives
\[
 \sigma_n=\sigma_{n-1}/P_{v_n}
                  (\gamma_n,\delta_n,\sigma_{n-1}),           \tag{6}
\]
so the PE current-coloring field is exact.  In SE the corresponding field
is simply \(\sigma_n=\sigma_{n-1}\).

Finally we transport the priority field excluding RE.  Suppose, contrary
to the required new field, that a
\noderef{RevisitingExtensionCertificate} existed under
\(\sigma_{n-1}\).  Retain its earlier PE index (h), exact ordered
colors \(\gamma_h,\delta_h\), individual boundary edge, inside and
outside endpoints, full cycle, and return path.  Each edge identity of
the return path is incident with (T_n), has one of these two colors,
and is controlled because (S_h) is contained in the accumulated raw
stored colors.  The symmetric stability criterion \noderef{sc1} therefore
copies that exact path, edge by edge, to \(\varphi_{n-1}\), retaining the
same connecting edge and endpoints.  If the containing old two-color
component were a path, the earlier PE exit path, whose delta is missing on
the retained tree, would be a distinct path component meeting (T_n).
When \(n>1\), lower assertion (ii) for the preceding ETT forbids two such
path components by \noderef{TreeInterchangeable}; when \(n=1\), the
positive earlier index required by the alleged certificate is impossible.
Hence in every possible case the copied component is a full cycle.
Together with the preserved return path and edge identity it gives an actual
\noderef{RevisitingExtensionCertificate} for the original data at slot
(n-1), contradicting the old SE or PE constructor's no-RE field.
Thus RE remains unavailable under \(\sigma_{n-1}\).

All fields of the same SE or PE constructor are now available.  In SE,
append the exact seed (T_n+f_n) and the composed closure from that seed
through (T^*) to (C).  In PE, append the composed closure from (T_n)
through (T^*) to (C).  The stored set, connecting-edge singleton,
chosen and outside vertices, gamma/delta fields, and extension kind are
the exact original fields, so parallel edge identities have never been
replaced by endpoint pairs.

We finish by assembling the certificate required for the arbitrary
(T^*).  Adjoin the rebuilt terminal constructor to the certified
replacement prefix.  This is a \noderef{TashkinovSeries} with (n)
rungs and final closure tree (C).  By (1), (T^*) strictly extends its
penultimate ladder tree (T_n) and is an ordered initial segment of (C),
so \noderef{ExtendedTashkinovTree} makes (T^*) its displayed replacement
ETT with generating coloring \(\sigma_n\).  The prefix construction gives
agreement of all ladder trees through (T_n), all raw (S_i), exact
(F_i), and kinds.  Its replacement colorings are proper and leave
exactly (e_0) uncolored; their positive stages are stable relative to
the old stages by the lower corresponding certificate, while the terminal
stage is stable by the hypothesis on \(\sigma_n\).  These are precisely
all fields of \noderef{CorrespondingExtendedTashkinovTree}; no equality
of the nondeterministic internal PE or RE witnesses is required by that
definition.

We have therefore constructed
\[
 \operatorname{Nonempty}
   (\noderef{CorrespondingExtendedTashkinovTree}
       (G,k,e_0,T,\sigma_n,T^*))
\]
for the arbitrary eligible ordered extension (T^*), using a closure
only to supply its final certified tuple.  This proves the universal clause
of \noderef{ModColoring}, and hence \(\sigma_n\) is a
\(\varphi_n\bmod T_n\) coloring as claimed.
\end{proof}
